\section{Related Work}
\label{sec:related-work}

\subsection{Manipulation learning and generalist policies}

Imitation learning addresses sequential distribution mismatch primarily during
training. DAgger collects expert labels on states induced by the learner and
adds them to the training set \cite{dagger}. This exposes the policy to recovery
states that are absent from the original demonstrations, but requires
interactive expert supervision. It does not provide a confidence signal for a
frozen policy at deployment.

Action-sequence policies address a different aspect of closed-loop behavior:
the temporal granularity of prediction and execution. ACT uses a conditional
variational autoencoder to generate action chunks and temporally ensembles
overlapping predictions \cite{act}. Diffusion Policy represents multimodal
action sequences through conditional denoising and executes them in a
receding-horizon loop \cite{diffusionpolicy}. Chunking moves the unit of
prediction from one action to a sequence and can smooth local behavior, while
making the whole proposed sequence the relevant unit of reliability.

Generalist VLA policies place semantic task interpretation and continuous
control within one learned pipeline \cite{kawaharazuka2025vla}. Their episode
outcomes consequently combine several capabilities, including visual
grounding, instruction interpretation and action generation. LIBERO evaluates
transfer across lifelong manipulation tasks, while LIBERO-PRO applies
controlled distribution shifts to the benchmark \cite{libero,liberopro}.
Task-level success measures the combined capability. Query-level analysis
provides a different resolution by following how policy evidence evolves as
successive action chunks alter the scene.

\subsection{World models and action-generating video models}

World-model policies introduce predicted consequences into action generation,
but differ in how prediction is connected to control. RoboDreamer treats a
generated video as an intermediate plan and recovers actions through inverse
dynamics \cite{robodreamer}. Unified World Models instead couples action and
future-observation diffusion, while Cosmos Policy jointly produces action
chunks, future observations and predicted returns
\cite{unifiedwm,cosmospolicy}. The first design separates visual planning from
action decoding; the latter designs place both outputs in one generative
process.

This distinction determines where uncertainty can enter. A staged video policy
can accumulate error in the predicted trajectory and in its action decoder,
whereas a joint model links action quality to a shared generative
representation. Both incur predictive computation that a direct VLA does not
require. At the same time, their predicted futures expose intermediate
structure that can be checked for consistency. Prior systems establish
video-conditioned action generation as a viable policy family. Pairing such a
policy with a separately trained direct policy additionally requires deciding
when its predictive process is likely to be useful.

\subsection{Uncertainty and runtime failure prediction}

That decision depends on uncertainty evidence whose interpretation matches the
generating process. Heteroscedastic likelihood learns input-dependent residual
variance and represents conditional, or aleatoric, dispersion under the fitted
model \cite{kendallgal}. Gaussian NLL induces inverse-variance weighting that
can progressively underweight poorly fitted regions; Beta-NLL modifies this
weighting to reduce the imbalance \cite{betanll}. Model uncertainty is commonly
approximated through disagreement among independently trained predictors
\cite{deepensembles}. For large robot policies, however, maintaining several
models adds training, memory and inference cost. Structured action outputs also
make a single global confidence value difficult to interpret: VLA studies
report that informative uncertainty varies across action dimensions and over
time \cite{criticaluq,vlacalibration}.

A related line of work converts such evidence into temporal failure scores.
FAIL-Detect learns a scalar score from successful trajectories and applies
time-varying calibrated thresholds \cite{faildetect}. FIPER requires only
successful rollouts: it aggregates observation-embedding novelty from RND-OE
and grid-based action-chunk entropy over short windows, calibrates both signals
and raises an alarm through their logical conjunction \cite{fiper}. SAFE instead
learns failure prediction from temporal VLA features and studies transfer to
held-out tasks \cite{safevla}. Our VLA monitor follows the latter supervised
setting but combines physical history, generated chunks, internal flow
uncertainty and candidate disagreement in one learned risk score.

Calibration also determines what can be claimed about a detector. Conformal
prediction provides finite-sample marginal coverage under exchangeability
\cite{conformal}; FIPER's bound concerns alarms on successful in-distribution
rollouts, not failure recall or coverage under arbitrary shift. Moreover, prior
VLA detectors monitor one policy distribution. They do not establish that
features or thresholds calibrated for a direct action model transfer to a
predictive policy with different outputs and error mechanisms.
